\section{Introduction}
\label{sec:intro}

Traffic forecasting underpins urban services---ride-hailing dispatch, signal
timing, congestion management---and STGNNs are the dominant
tool~\citep{li2018dcrnn,wu2019gwnet,yu2018stgcn}. In the lab, these models are
evaluated on held-out data from the same city, sensors, and time window as
training. Deployment conditions differ: sensors fail, models trained in one
city are applied to another, and traffic shifts between weekdays and weekends.
Which architectures hold up under these shifts is a deployment question that
in-distribution accuracy cannot answer.

Prior STGNN robustness work emphasizes synthetic
corruptions~\citep{zhang2025strap} or frames cross-city transfer as few-shot
learning~\citep{zeng2025fedtt,jing2026tlgpstgn}. Missing is a controlled
\emph{diagnostic} study that (i) uses deployment-realistic shifts, (ii)
evaluates the same frozen backbone under multiple shifts, (iii) attributes
differences to specific architectural choices, and (iv) pre-registers
hypotheses so refuted predictions carry the same evidential weight as
confirmed ones. Cross-city comparisons also conflate transfer difficulty with
inductive bias---adaptive models can re-learn target topology, fixed models
cannot---which a matched-structural-access protocol (source-city topology for
all backbones) removes. No prior study systematically compares how different
adaptive-relation designs---GWNet's dense softmax routing, MegaCRN's
meta-graph memory, STAEformer's adaptive embeddings---redistribute fragility
across deployment shifts.

We assemble a bank of eight architectures trained on METR-LA---a five-model
canonical bank spanning graph-free (GRU), fixed-graph (DCRNN, STGCN), and
adaptive-graph (MTGNNLight, GWNet), plus three modern adaptive-relation models
(AGCRN, MegaCRN, STAEformer)---and evaluate each frozen checkpoint under three
shifts: \textbf{sensor dropout} (10--50\% zeroed at inference, 5 random masks),
\textbf{cross-city transfer} (zero-shot METR-LA $\to$ PEMS-BAY under
matched structural access), and \textbf{temporal shift} (weekend vs.\ weekday
windows within the METR-LA test set).

Three findings stand out:
\begin{enumerate}
  \item No architecture dominates all three shifts; rankings change
        qualitatively between shift types. GWNet's dense adaptive adjacency
        is most fragile under sensor dropout (${\sim}4.3{\times}$ worse than
        DCRNN at 30\%) but not the worst at city transfer. DCRNN leads
        cross-city in the canonical bank; MegaCRN overtakes it once the
        extended set is added.

  \item GWNet's dropout fragility is the visible tail of a routing-density
        relation: $\mathrm{max\_mass}(A)$ ranks $\Delta_{\mathrm{sensor}}$ at
        Pearson $\rho{=}{+}0.979$ over four adaptive backbones on METR-LA,
        replicating at $\rho{=}{+}0.910$ on PEMS04 with native checkpoints.
        Replacing the learned adjacency with a fixed road graph cuts dropout
        sensitivity by 46\%; a mask-aware variant recovers 75\% without
        retraining. AGCRN ($\mathrm{max\_mass}{=}0.008$) and STAEformer match
        DCRNN despite learning adaptive relations: the operative factor is
        routing density, not topology learning.

  \item Modern adaptive-relation models distribute their failures across axes.
        AGCRN and STAEformer survive sensor dropout; MegaCRN leads cross-city.
        But MegaCRN and STAEformer carry weekday/weekend gaps
        ${\sim}8$--$9{\times}$ the canonical bank---an axis that the
        five-model evaluation could not have surfaced, and one that explicit
        TOD/DOW covariates do not explain
        (Section~\ref{sec:mechanism}).
\end{enumerate}

Sixteen predictions were pre-registered in five stages, each frozen before
the next round of training (ledger and confirmed/refuted/partial criteria:
Appendix~\ref{sec:appendix_preregistered}). Four were refuted: GWNet's
dominant failure axis was predicted to be cross-city transfer, not sensor
dropout; temporal fragility was predicted to be architecture-invariant, but
is not---findings invisible in a confirmation-only evaluation. Three
follow-ups---PEMS04-native routing-density replication
(Appendix~\ref{sec:appendix_pems04_routing}), the mask-aware dropout-ratio
sweep (Appendix~\ref{sec:appendix_mask_aware_sweep}), and cross-city
zero-padding sensitivity (Appendix~\ref{sec:appendix_no_pad_sensitivity})---are
flagged as reviewer-driven, reported alongside the pre-registered core. The
frozen-checkpoint protocol is reproducible end-to-end on a single commodity
GPU; checkpoints are released as supplementary material.
